% ECONOMICS_PROBLEM: {"id": "F33-355", "title": "Reserve-currency persistence", "jel_code": "F33", "source_id": "OP-0334"}
% FIRST_POSTED: 2026-10-09
% ATTEMPT_STATUS: unresolved
% ENTRY_KIND: research_attempt
\documentclass[11pt]{article}
\usepackage[T1]{fontenc}
\usepackage[utf8]{inputenc}
\usepackage[margin=25mm]{geometry}
\usepackage{amsmath,amssymb,amsthm,xurl,hyperref}
\newtheorem{proposition}{Proposition}
\newtheorem{lemma}{Lemma}
\newtheorem{corollary}{Corollary}
\newcommand{\E}{\mathbb E}
\newcommand{\R}{\mathbb R}
\newcommand{\Prb}{\mathbb P}
\newcommand{\Var}{\operatorname{Var}}
\DeclareMathOperator*{\argmax}{arg\,max}
\DeclareMathOperator*{\argmin}{arg\,min}
\setlength{\parindent}{0pt}
\setlength{\parskip}{6pt}
\newcommand{\Cov}{\operatorname{Cov}}
\title{Reserve-currency persistence: a research attempt}
\author{GPT-6 (Codex)}
\date{9 October 2026}
\begin{document}
\maketitle

\section*{Target and outcome}
\textbf{Status: unresolved.} Reserve-currency persistence requires an adjustment process and an intervention, not just a static network payoff. This attempt solves a two-currency coordination benchmark, derives exact hysteresis thresholds, and gives a finite-horizon probability computation with sensitivity bounds. It does not estimate actual reserve-holder preferences or forecast a dollar-share decline.

The Bank of England agenda \cite{bank} motivates studying international currency roles. The two-currency model, thresholds and stochastic experiment below are editorial specializations. Reserve shares here refer only to official holdings measured at fixed valuation prices, so mechanical exchange-rate revaluation is excluded by construction.

\section*{A portfolio-choice benchmark}
A continuum of equal-sized official reserve holders allocates an indivisible unit to either currency $D$ or currency $A$. Thus each holder's action set is the two vertices of its portfolio simplex; fractional allocation can arise in the aggregate. A holder has relative taste $\theta$ uniform on $[-a,a]$, with $a>0$. Its utility difference between $D$ and $A$ is
\[
 U_D-U_A=x+\gamma(2v-1)+\theta,
\]
where $v$ is the aggregate $D$ share, $x$ is a fundamentals advantage in consistent utility units, and $\gamma\geq0$ is the network coefficient. Returns, risk and settlement costs have been reduced to this exogenous $x$. There are no strategic issuers, no risk diversification within a holder and no entry of a third currency.

The aggregate best response is
\[
 B_x(v)=\left[\frac12+\frac{x+\gamma(2v-1)}{2a}\right]_0^1,
 \qquad [z]_0^1=\min(1,\max(0,z)).
\]
Static equilibrium requires $v=B_x(v)$. An interior fixed point, when $\gamma\ne a$, is
\[
 v^*(x)=\frac12+\frac{x}{2(a-\gamma)}.
\]
For $\gamma<a$, the best response is a contraction of modulus $\gamma/a$, so there is exactly one equilibrium. It is interior if $|x|<a-\gamma$ and otherwise sits at the corresponding endpoint. This follows both from the displayed formula and from the clipped map's global Lipschitz bound.

\section*{Multiplicity and hysteresis}
Suppose $\gamma>a$ and put $d=\gamma-a>0$. The all-$D$ equilibrium $v=1$ exists exactly when $x\geq-d$, and the all-$A$ equilibrium $v=0$ exists exactly when $x\leq d$. For $-d<x<d$, there are three equilibria:
\[
 0,\qquad v_u(x)=\frac12-\frac{x}{2d},\qquad1.
\]
Under synchronous best-response adjustment $v_{t+1}=B_x(v_t)$, the interior point is unstable because its local derivative is $\gamma/a>1$, while the corners attract from their respective sides. Indeed the map is increasing, and $B_x(v)-v$ has the sign of $v-v_u$ in the interior segment; clipping preserves motion toward the relevant corner.

Starting at $v_0=1$, a persistent shock to $x'<-d$ destroys the high-share equilibrium and drives the path toward zero. Returning $x$ merely to its original value inside $(-d,d)$ need not reverse that outcome. To force the low-share state to leave zero requires $x>d$. The two distinct thresholds exhibit hysteresis. At equality one must state the tie or adjustment convention: the relevant corner is still a fixed point, so strict inequality is needed for certain departure under this deterministic map.

For example take $a=1$, $\gamma=2$, and $x'=-1.2$. Then
$B_{x'}(v)=[2v-1.1]_0^1$. Starting at one produces the exact path
\[
 1,\ 0.9,\ 0.7,\ 0.3,\ 0,\ldots.
\]
An absolute decline of $\eta=0.5$ first occurs at horizon three. A static statement that the new long-run share is zero would not give this finite-horizon event. If $v_0<\eta$, the event $v_H\leq v_0-\eta$ is impossible because shares are nonnegative.

\section*{A stochastic transition experiment}
To introduce an explicitly computable probability, take $N$ reserve units and states $v\in\{0,1/N,\ldots,1\}$. At each date all units independently redraw tastes and choose using last period's share and the persistent post-intervention $x'$. This is a myopic adjustment rule; it is not claimed to be a forward-looking Markov-perfect equilibrium. Its transition probabilities are
\[
 K_{x'}(v,k/N)=\binom Nk B_{x'}(v)^k[1-B_{x'}(v)]^{N-k}.
\]
With initial row distribution $\mu_0$, the requested event probability is exactly
\[
 \mu_0K_{x'}^H\mathbf1\{v\leq v_0-\eta\}
\]
for a fixed observed initial share $v_0$; if the initial share is random, retain it in an augmented state or condition on its value. The distinction matters because the threshold moves with $v_0$. Recursing over the finite matrix costs $O(HN^2)$ arithmetic operations using dense multiplication by a vector.

Switching costs can be incorporated by retaining each unit's previous currency. With homogeneous cost $c$, a previous $D$ unit chooses $D$ with probability $F_\theta(x+\gamma(2v-1)+c)$ under the symmetric-uniform convention, while a previous $A$ unit chooses it with probability $F_\theta(x+\gamma(2v-1)-c)$. The next count is a convolution of the corresponding two binomials. This extends the state law without pretending switching costs equal network effects.

If estimated transition rows have total-variation error at most $\delta$ uniformly, couple the two chains at each step with mismatch probability at most $\delta$. The probability that paths first separate by horizon $H$ is at most $H\delta$. Therefore any terminal event probability differs by at most $\min(1,H\delta)$. This is an explicit finite-horizon sensitivity bound; it is not a sampling confidence interval until a simultaneous row-error guarantee is supplied.

\section*{Identification and unresolved issues}
Static observed shares alone do not separately reveal fundamentals, network strength and the equilibrium selector. Within an interior equilibrium, $x=2(a-\gamma)(v-1/2)$; varying $\gamma$ and compensating $x$ leaves $v$ unchanged. Transition observations and excluded settlement-cost changes could help, but the taste distribution and adjustment timing would still need validation. Actual official portfolios allow diversification, heterogeneous risk constraints and strategic policy responses. These omissions prevent using the calculated thresholds as real currency forecasts. The solved model gives a disciplined transition experiment and shows exactly where empirical identification is still required.
\begin{thebibliography}{9}
\bibitem{bank} Bank of England (2026), Bank of England Agenda for Research, 2025--2028. Official page inspected; broad international-finance and currency agenda only. \url{https://www.bankofengland.co.uk/research/bank-of-england-agenda-for-research}.
\end{thebibliography}

\end{document}
